\documentclass[aspectratio=169]{beamer}

\usepackage{amsmath,amssymb}
\usepackage{listings}

\usepackage{style}
%%% Author, Affiliation, Date
\author{Sheng-Chieh Tseng}
\institute{Department of Physics, National Taiwan University}
\date{5 Jun 2026}



%%% defs
\def\cmd#1{\texttt{\color{red}\footnotesize $\backslash$#1}}
\def\env#1{\texttt{\color{blue}\footnotesize #1}}
%\definecolor{erasmusmc}{RGB}{134,210,236}
%\definecolor{erasmusmc_dark}{RGB}{12,33,116}
\definecolor{deepblue}{rgb}{0,0,0.5}
\definecolor{deepred}{rgb}{0.6,0,0}
\definecolor{deepgreen}{rgb}{0,0.5,0}
\definecolor{halfgray}{gray}{0.55}


% Rounded blocks with shadow for a consistent report style.
\setbeamertemplate{blocks}[rounded][shadow=true]

\lstset{
    basicstyle=\ttfamily\small,
    keywordstyle=\bfseries\color{erasmusmc_dark},
    emphstyle=\ttfamily\color{erasmusmc!50!erasmusmc_dark},    % Custom highlighting style
    stringstyle=\color{erasmusmc!50!deepgreen},
    numbers=left,
    numberstyle=\small\color{halfgray},
    rulesepcolor=\color{red!20!green!20!blue!20},
    frame=shadowbox,
}

\renewcommand{\familydefault}{\sfdefault} % You can define your font style here


% Resolve figures whether the deck is compiled from report/ or repo root.
\graphicspath{%
    {../results/report/result/cuda_mg_phi_plots/plots/}%
    {../results/report/result/cuda_mg_convergence/}%
    {../results/report/result/omp_residual_history/}%
    {../results/report/result/omp_thread_scaling/}%
    {results/report/result/cuda_mg_phi_plots/plots/}%
    {results/report/result/cuda_mg_convergence/}%
    {results/report/result/omp_residual_history/}%
    {results/report/result/omp_thread_scaling/}%
}

\title[CUDA Poisson MG 2D/3D]{CUDA-Accelerated Multigrid Poisson Solvers in 2D and 3D}

\begin{document}

\begingroup
\AtBeginSection[]{}
\section{Problem Setup}
\setbeamertemplate{title page}{
    \vbox to \textheight{
        \vspace*{2cm}
        \begin{minipage}[t][0.275\textheight][t]{\textwidth}
            \begin{center}
                {\usebeamercolor[fg]{normal text}\bfseries\resizebox{0.95\textwidth}{!}{\inserttitle}\par}
            \end{center}
        \end{minipage}\par

        \begin{minipage}[c][0.15\textheight][c]{\textwidth}
            \begin{center}
                {\usebeamerfont{author}\insertauthor\par}
                \vspace{0.5em}
                {\usebeamerfont{institute}\insertinstitute\par}
                \vspace{0.5em}
                {\usebeamerfont{date}\insertdate\par}
            \end{center}
        \end{minipage}\par

        \begin{minipage}[b][0.25\textheight][b]{\textwidth}
            \begin{center}
                \begin{minipage}{0.92\textwidth}
                    \begin{columns}[c,onlytextwidth]
                        \column{0.76\textwidth}
                        \raggedright
                        {\small\textbf{Project repository:}\par}
                        \vspace{0.12em}
                        {\footnotesize\url{https://github.com/TCC0731/multigrid-poisson-solvers}}
                        \column{0.24\textwidth}
                        \centering
                        \raisebox{-0.08cm}{
\includegraphics[width=1.9cm]{qrcode-generator.png}}
                    \end{columns}
                \end{minipage}
            \end{center}
        \end{minipage}

        \vfill
    }
}

\begin{frame}[plain]
    \titlepage
\end{frame}

\begin{frame}
    \tableofcontents[sectionstyle=show,subsectionstyle=show/shaded/hide,subsubsectionstyle=show/shaded/hide]

\end{frame}

\endgroup

\begin{frame}[t]{Problem definition}
    \footnotesize
    \begin{block}{Continuous model}
        We solve the positive-definite Poisson problem with Dirichlet boundary conditions:
        \[
            \mathcal{L}\phi = f, \qquad \mathcal{L} = -\nabla^2,
        \]
        on \(\Omega = [0,1]^d\), with \(d \in \{2,3\}\).
    \end{block}

    \begin{block}{Discrete model}
        \[
            h = \frac{1}{n+1}, \qquad \phi_0\vert_{\text{interior}} = 0, \qquad
            \phi\vert_{\partial\Omega} = \phi_{\text{exact}}\vert_{\partial\Omega}.
        \]
        The initial guess is zero in the interior, while boundary values are fixed by the manufactured solution.
        The code uses the standard 5-point stencil in 2D and 7-point stencil in 3D.
    \end{block}
\end{frame}

\begin{frame}[t]{Problem definition (cont.)}
    \footnotesize
    \begin{block}{Dimensionless discrete operator}
        The interior unknowns satisfy
        \[
            A_h \phi_h = h^2 f_h, \qquad r_h = h^2 f_h - A_h \phi_h.
        \]
        In 2D, the standard 5-point stencil is
        \[
            (A_h \phi)_{i,j} = 4\phi_{i,j} - \phi_{i-1,j} - \phi_{i+1,j} - \phi_{i,j-1} - \phi_{i,j+1},
        \]
        and in 3D the 7-point stencil replaces \(4\) by \(6\) and includes the six axis-aligned neighbors.
    \end{block}

    \begin{itemize}
        \item Dirichlet boundary data are taken from the manufactured solution; sine gives homogeneous Dirichlet data, while cosine gives inhomogeneous Dirichlet data.
        \item The arrays store \(n\) interior points per axis plus one boundary layer on each side, so the total size is \(n+2\).
    \end{itemize}
\end{frame}

\begin{frame}[t]{Solver conventions}
    \footnotesize
    \begin{itemize}
        \item Convergence criterion: the relative scaled residual
              \[
                  \eta = \frac{\|h^2 f - A_h\phi\|_2}{\|h^2 f\|_2}
              \]
              where \(A_h\) is the dimensionless finite-difference stencil approximating \(h^2(-\nabla^2)\),
              must satisfy \(\eta \le \texttt{tol}\).
        \item Baseline smoothers: red-black SOR.
        \item Multigrid smoothing: red-black SOR with \(\nu\) pre- and post-smoothing steps.
        \item If \(\omega = 1\), the red-black SOR update reduces to red-black Gauss--Seidel.
    \end{itemize}

    \begin{block}{Automatic SOR weight}
        When \(\omega\) is left on auto, the code uses
        \[
            \omega_{\mathrm{auto}} = \frac{2}{1 + \sin\!\bigl(\pi/(n+1)\bigr)}.
        \]
        Here \(n\) is the number of interior points per axis on the current grid level.
    \end{block}
\end{frame}

\begin{frame}[t]{Solver overview}
    \footnotesize
    \begin{columns}[T,onlytextwidth]
        \column{0.48\textwidth}
        \begin{block}{Implemented solvers}
            \begin{itemize}
                \item Red-black SOR baseline.
                \item Multigrid Poisson solver in 2D and 3D.
                \item Both V-cycle and W-cycle variants.
                \item Coarsest-grid solve with SOR or exact/direct solve.
            \end{itemize}
        \end{block}

        \column{0.48\textwidth}
        \begin{block}{Implemented support}
            \begin{itemize}
                \item Python backend.
                \item OpenMP CPU backend.
                \item CUDA GPU backend.
                \item Full-weighting restriction and bilinear/trilinear prolongation.
                \item Manufactured-solution Dirichlet boundary setup.
            \end{itemize}
        \end{block}
    \end{columns}
\end{frame}

\begin{frame}[t]{Multigrid operators}
    \footnotesize
    \begin{itemize}
        \item Restriction: \(r_{2h} = R\,r_h\), where the 2D full-weighting stencil is
              \[
                  R_{2D} = \frac{1}{16}
                  \begin{bmatrix}
                      1 & 2 & 1 \\
                      2 & 4 & 2 \\
                      1 & 2 & 1
                  \end{bmatrix},
              \]
              and the 3D version is the corresponding tensor-product stencil
              \[
                  (Rr)_{I,J,K} = \frac{1}{64}\sum_{a,b,c=-1}^{1} w_a w_b w_c\, r_{2I+a,\,2J+b,\,2K+c},
                  \qquad w_0 = 2,\; w_{\pm1} = 1.
              \]
        \item Prolongation: \(e_h = P\,e_{2h}\), where 2D uses bilinear interpolation and 3D uses trilinear interpolation; for example,
              \[
                  e_h(2I+1, 2J) = \frac{1}{2}\bigl(e_{2h}(I, J) + e_{2h}(I+1, J)\bigr),
                  \quad
                  e_h(2I+1, 2J+1) = \frac{1}{4}\sum_{p,q \in \{0,1\}} e_{2h}(I+p, J+q),
              \]
              and the 3D case uses the analogous tensor-product trilinear interpolation, involving 2, 4, or 8 coarse-grid neighbors depending on the fine-grid point location.
        \item Coarsest-grid solve: exact solve by default, with an RB-SOR coarse-solve variant also available.
    \end{itemize}
\end{frame}

\begin{frame}[t]{Coarsest-grid Solver Design}
    \footnotesize
    \begin{block}{Why the bottom level matters}
        \begin{itemize}
            \item The coarsest solve closes the multigrid hierarchy.
            \item A weak bottom-level solve can leak error back to the fine grid after prolongation.
            \item The choice affects both robustness and time-to-solution.
        \end{itemize}
    \end{block}

    \begin{columns}[T,onlytextwidth]
        \column{0.48\textwidth}
        \begin{block}{Our implementation}
            \begin{itemize}
                \item Default: exact/direct solve on the coarsest grid.
                \item Optional: RB-SOR coarse solve for comparison or fallback.
            \end{itemize}
        \end{block}

        \column{0.48\textwidth}
        \begin{block}{Trade-off}
            \begin{itemize}
                \item RB-SOR keeps the code path simple and is easy to benchmark.
                \item Exact/direct gives the strongest bottom-level correction when the coarse system is tiny.
            \end{itemize}
        \end{block}
    \end{columns}
\end{frame}

\begin{frame}[t]{Multigrid workflow}
    \footnotesize
    \begin{block}{One V-cycle}
        \begin{itemize}
            \item Pre-smooth on the current grid to damp high-frequency error.
            \item Compute the residual \(r_h = h^2 f - A_h \phi\), where \(A_h\) is the dimensionless discrete Poisson stencil.
            \item Restrict the residual to the coarse grid: \(r_{2h} = R\,r_h\).
            \item Solve the coarse-grid error equation recursively, \(A_{2h} e_{2h} = r_{2h}\).
            \item Prolongate the correction and update the solution: \(\phi \leftarrow \phi + P\,e_{2h}\).
            \item Post-smooth and repeat until the relative scaled residual satisfies \(\eta \le \texttt{tol}\).
        \end{itemize}
    \end{block}

    \[
        \begin{aligned}
            \text{pre-smooth} & \rightarrow \text{compute residual }(r_h) \rightarrow \text{restrict }(r_{2h}=Rr_h)                                                                         \\
                              & \rightarrow \text{coarse solve }(A_{2h}e_{2h}=r_{2h}) \rightarrow \text{prolongate/update }(\phi \leftarrow \phi + P e_{2h}) \rightarrow \text{post-smooth}
        \end{aligned}
    \]
\end{frame}

\begin{frame}[t]{W-cycle and MG grid sizes}
    \footnotesize
    \begin{block}{W-cycle}
        \begin{itemize}
            \item Start like a V-cycle: pre-smooth, form the residual, and restrict to the next coarser grid.
            \item Call the coarse level twice before coming back up to the fine grid.
            \item The repeated coarse visit reduces the remaining low-frequency error, but it also increases the per-cycle cost.
        \end{itemize}
    \end{block}

    \[
        \begin{aligned}
            \text{pre-smooth} & \rightarrow \text{compute residual }(r_h) \rightarrow \text{restrict }(r_{2h}=Rr_h)                                                            \\
                              & \rightarrow \text{coarse cycle} \times 2 \rightarrow \text{prolongate/update }(\phi \leftarrow \phi + P e_{2h}) \rightarrow \text{post-smooth}
        \end{aligned}
    \]

    \begin{block}{MG-compatible grid sizes}
        \begin{itemize}
            \item \texttt{--grid-size} is the number of interior points per axis, so the full array size is \texttt{grid-size + 2}.
            \item Nested coarsening uses \(n_c = (n_f - 1)/2\), and the recursion stops when \(n \le 4\).
            \item Valid families: \(2^a - 1\), \(3 \cdot 2^a - 1\), and \(5 \cdot 2^a - 1\).
        \end{itemize}
    \end{block}
\end{frame}

\begin{frame}[t]{Time cost}
    \footnotesize
    \setlength{\abovedisplayskip}{4pt}
    \setlength{\belowdisplayskip}{4pt}
    \setlength{\abovedisplayshortskip}{2pt}
    \setlength{\belowdisplayshortskip}{2pt}
    \begin{block}{Coarse-grid scaling}
        Assume the fine-grid work is \(N\). One level down:
        \[
            2D: N \rightarrow N/4, \qquad 3D: N \rightarrow N/8
        \]
    \end{block}

    \begin{columns}[T,onlytextwidth]
        \column{0.475\textwidth}
        \begin{block}{V-cycle}
            \[
                \begin{aligned}
                    T_V      & \sim N + \frac{N}{2^d} + \frac{N}{2^{2d}} + \cdots   \\
                    T_V^{2D} & \sim \frac{4}{3}N, \qquad T_V^{3D} \sim \frac{8}{7}N
                \end{aligned}
            \]
        \end{block}

        \column{0.475\textwidth}
        \begin{block}{W-cycle}
            \[
                \begin{aligned}
                    T_W      & \sim N + 2\frac{N}{2^d} + 4\frac{N}{2^{2d}} + \cdots \\
                    T_W^{2D} & \sim 2N, \qquad T_W^{3D} \sim \frac{4}{3}N
                \end{aligned}
            \]
        \end{block}
    \end{columns}

    \begin{block}{W / V ratio and takeaway}
        \(2D: T_W/T_V = 3/2 \approx 50\%, \quad 3D: T_W/T_V = 7/6 \approx 17\%\).
        3D shrinks faster, so the extra W-cycle coarse-grid work matters less than in 2D.
    \end{block}
\end{frame}

\begin{frame}[t]{Mixed sine result: 2D}
    \centering
    \includegraphics[width=0.98\textwidth,height=0.74\textheight,keepaspectratio]{phi_2d_v_sor_mixed_sine_n4095_nu3_w1p25_cs16.png}

    \vspace{0.35em}
    {\footnotesize CUDA MG V-SOR on the manufactured solution \(\phi_{\rm exact}(x,y)=\sin(2\pi x)\sin(3\pi y)\) with \(\nu = 3\) and \(\omega = 1.25\).}
\end{frame}

\begin{frame}[t]{Mixed sine result: 3D}
    \centering
    \includegraphics[width=0.98\textwidth,height=0.74\textheight,keepaspectratio]{phi_3d_v_sor_mixed_sine_n383_nu3_w1p25_cs16.png}

    \vspace{0.35em}
    {\footnotesize CUDA MG V-SOR on the manufactured solution \(\phi_{\rm exact}(x,y,z)=\sin(2\pi x)\sin(3\pi y)\sin(4\pi z)\) with \(\nu = 3\) and \(\omega = 1.25\).}
\end{frame}

\begin{frame}[t]{CUDA MG Convergence: L2 Error}
    \centering
    \includegraphics[width=0.96\textwidth,height=0.80\textheight,keepaspectratio]{plots_2d_3d_l2_error_sine.png}
\end{frame}

\begin{frame}[t]{CUDA MG Convergence: SOR and Exact Iteration Time}
    \centering
    \includegraphics[width=0.96\textwidth,height=0.80\textheight,keepaspectratio]{plots_2d_3d_iter_time_sine.png}
\end{frame}

\begin{frame}[t]{CUDA MG Convergence: Iteration Time}
    \centering
    \includegraphics[width=0.96\textwidth,height=0.80\textheight,keepaspectratio]{plots_2d_3d_iter_time_fit_last10_sine.png}
\end{frame}

\begin{frame}[t]{OpenMP residual history}
    \centering
    \includegraphics[width=0.96\textwidth,height=0.8\textheight,keepaspectratio]{residual_history.png}
\end{frame}

\begin{frame}[t]{OpenMP thread scaling}
    \centering
    \includegraphics[width=0.96\textwidth,height=0.8\textheight,keepaspectratio]{plots_2d_3d.png}
\end{frame}

\end{document}
